\documentclass[11pt]{article}

\usepackage[margin=1in]{geometry}
\usepackage{amsmath,amssymb,amsthm,mathtools}
\usepackage{booktabs,tabularx,array,longtable}
\usepackage[section]{placeins}
\usepackage{float}
\usepackage{microtype}
\usepackage{graphicx}
\usepackage[numbers,sort&compress]{natbib}
\usepackage[colorlinks=true,allcolors=blue]{hyperref}
\usepackage[nameinlink]{cleveref}
\usepackage[T1]{fontenc}
\usepackage[utf8]{inputenc}
\usepackage{xcolor}
\usepackage{listings}
\usepackage{listingsutf8}
\usepackage{authblk}
\setlength{\emergencystretch}{3em}

% ArXiv-friendly Lean code: no shell escape, minted, or external highlighter.
\definecolor{leankeyword}{HTML}{243B6B}
\definecolor{leantactic}{HTML}{70477A}
\definecolor{leantype}{HTML}{1F6473}
\definecolor{leancomment}{HTML}{39734D}
\definecolor{leanstring}{HTML}{8A4B2A}
\definecolor{leannumber}{HTML}{777777}
\definecolor{leanbackground}{HTML}{F7F7F3}
\lstdefinelanguage{lean}{%
  sensitive=true,
  morecomment=[l]{--},
  morecomment=[s]{/-}{-/},
  morestring=[b]",
  morekeywords=[1]{import,open,namespace,section,end,noncomputable,variable,variables,
    universe,universes,set_option,attribute,axiom,opaque,theorem,lemma,def,abbrev,
    example,instance,structure,class,inductive,notation,local,deriving,extends,
    by,where,let,have,show,fun,match,with,if,then,else,do,return,in,forall},
  morekeywords=[2]{simp,simpa,rw,erw,dsimp,unfold,ring,ring_nf,nlinarith,linarith,
    positivity,norm_num,field_simp,omega,exact,refine,apply,assumption,intro,intros,
    ext,funext,cases,rcases,obtain,constructor,left,right,classical,by_contra,
    contrapose,convert,using,suffices,calc},
  morekeywords=[3]{Prop,Type,Nat,Int,Real,Fin,Set,Submodule,Matrix,EuclideanSpace,
    LinearMap,ContinuousLinearMap,LinearPMap,RCLike,IsTrialResidual,
    IsExactSpectralDecomposition,FormBoundedSylvesterGap,
    NormalizedSymmetricOperatorIdealFamily,SourcePopulationGap},
}
\lstdefinestyle{leanpaper}{%
  language=lean,
  columns=fullflexible,
  backgroundcolor=\color{leanbackground},
  basicstyle=\ttfamily\scriptsize,
  commentstyle=\color{leancomment}\itshape,
  keywordstyle=[1]\color{leankeyword}\bfseries,
  keywordstyle=[2]\color{leantactic},
  keywordstyle=[3]\color{leantype},
  stringstyle=\color{leanstring},
  numberstyle=\tiny\color{leannumber},
  breakatwhitespace=false,
  breaklines=true,
  captionpos=b,
  keepspaces=true,
  numbers=left,
  numbersep=6pt,
  showspaces=false,
  showstringspaces=false,
  showtabs=false,
  tabsize=2,
  aboveskip=0.8em,
  belowskip=0.8em,
  literate=
    {\\KField}{{$\Bbbk$}}1
}
\crefname{lstlisting}{formalization}{formalizations}
\Crefname{lstlisting}{Formalization}{Formalizations}
\renewcommand{\lstlistingname}{Formalization}

\newtheorem{theorem}{Theorem}
\newtheorem{proposition}[theorem]{Proposition}
\newtheorem{remark}[theorem]{Remark}

% Generated by scripts/build_source_census.py; do not edit by hand.
\newcommand{\YWSCensusEntryCount}{24}
\newcommand{\YWSCensusPaperSourceCount}{22}
\newcommand{\YWSCensusExplicitAdditionCount}{2}
\newcommand{\YWSCensusVerifiedCount}{24}
\newcommand{\YWSCensusCorrectedCount}{7}

% Generated by scripts/fetch_openalex_bibliometrics.py; do not edit by hand.
\newcommand{\OpenAlexRetrievedAtUTC}{2026-08-17T20:33:55Z}
\newcommand{\OpenAlexRetrievedDate}{17 August 2026}
\newcommand{\OpenAlexCurrentYear}{2026}
\newcommand{\OpenAlexLatestCompleteYear}{2025}
\newcommand{\OpenAlexDavisKahanCitations}{1228}
\newcommand{\OpenAlexDavisKahanWorkID}{W1970377488}
\newcommand{\OpenAlexDavisKahanUpdatedDate}{2026-08-16T07:02:28.622633}
\newcommand{\OpenAlexDavisKahanCurrentYearCitations}{57}
\newcommand{\OpenAlexDavisKahanLatestCompleteYearCitations}{60}
\newcommand{\OpenAlexYuWangSamworthCitations}{462}
\newcommand{\OpenAlexYuWangSamworthWorkID}{W2088911135}
\newcommand{\OpenAlexYuWangSamworthUpdatedDate}{2026-08-16T07:02:28.622633}
\newcommand{\OpenAlexYuWangSamworthCurrentYearCitations}{22}
\newcommand{\OpenAlexYuWangSamworthLatestCompleteYearCitations}{49}


% Manually populated from Google Scholar on the same date as the OpenAlex observation.
\newcommand{\GoogleScholarDavisKahanCitations}{2086}
\newcommand{\GoogleScholarYuWangSamworthCitations}{1067}

\newcommand{\Lean}{\textsc{Lean}}
\newcommand{\YWS}{Yu--Wang--Samworth}
\newcommand{\DK}{Davis--Kahan}
\newcommand{\TauCeti}{Tau Ceti}
\newcommand{\DKResultCount}{29}
\newcommand{\DKProvedCount}{28}
\newcommand{\DKRefutedCount}{1}
\newcommand{\DKSemanticAcceptedCount}{29}

\title{Formalizing Davis--Kahan Perturbation Theory in Lean}
% Additional authors are expected; keep the confirmed ordering below stable.

\author[1]{Jonathan Crall}
\author[2]{Edward Wang}
\author[1]{Brian Hu}
\author[2]{Carey E. Priebe}

\affil[1]{%
  Kitware, Inc.\\
  \texttt{jon.crall@kitware.com},
  \texttt{brian.hu@kitware.com}
}

\affil[2]{%
  Johns Hopkins University\\
  \texttt{ewang43@jhu.edu},
  \texttt{cep@jhu.edu}
}

\date{Draft -- September 2026}

\begin{document}
\maketitle

\begin{abstract}
We formalize in Lean~4 the result-level mathematical content of Davis and
Kahan's 1970 paper on rotations of eigenspaces.  A source-based coverage rule
selects \DKResultCount{} targets: the four unnumbered theorems of Section~2 and
every named theorem, proposition, lemma, and corollary established in the
paper.  Twenty-eight are proved at the scope stated in the source.  The
remaining target, Proposition~4.4, is false as printed; we give a
machine-checked counterexample in $\mathbb{R}^4$, identify the failing step in
the published proof, and prove a corrected result for the $Q$-norm class.

The formalization retains the parts of the source that are easy to lose in a
finite-dimensional restatement: real and complex Hilbert spaces,
infinite-dimensional operator ideals, unbounded self-adjoint operators, domain
conditions, and the half-infinite spectral-separation cases used by the
single-angle sine theorem.  We use that theorem as a detailed correspondence
case study.  In particular, the positive Davis--Kahan $\sin\Theta_0$ operator
and the rectangular Lean representative $(I-F_0F_0^*)E_0$ act on different
spaces.  The modulus of the rectangular operator is the positive
$\sin\Theta_0$ operator, and polar decomposition therefore gives equality of
their supported unitarily invariant norms.

The proof required reusable formal theory for principal angles, approximation
numbers and symmetric operator ideals, Borel spectral calculus, partial
unbounded operators, real/complex transfer, direct rotations, and Sylvester
equations.  As a finite-dimensional statistical application, we also formalize
the theorem-level claims of Yu, Wang, and Samworth and record two printed source
defects found during that comparison.
\end{abstract}

\section{Introduction}
\label{sec:introduction}

Davis and Kahan's perturbation theory for spectral subspaces is a standard
bridge from operator perturbation bounds to geometric control of eigenspaces
\citep{davis1970rotation,kato1976perturbation,stewart1990matrix}.  Its most
recognizable consequence is the sine-theta estimate, but the 1970 paper develops
considerably more: tangent and double-angle bounds, canonical direct rotations,
majorization statements for unitarily invariant norms, Sylvester-operator
estimates, and extensions to unbounded self-adjoint operators.  These results
remain part of the foundation for modern matrix perturbation theory and
statistical spectral analysis \citep{wedin1972perturbation,chen2021spectral}.
On \OpenAlexRetrievedDate{}, OpenAlex indexed
\OpenAlexDavisKahanCitations{} citing works for Davis--Kahan and
\OpenAlexYuWangSamworthCitations{} for YWS; Google Scholar reported
\GoogleScholarDavisKahanCitations{} and
\GoogleScholarYuWangSamworthCitations{} citations, respectively.

Lean~4 is an interactive theorem prover whose kernel checks elaborated proof
terms, and Mathlib is its community mathematical library
\citep{demoura2021lean4,mathlib2020}.  In this project, formalization has two
separate obligations.  Kernel verification establishes that the theorem encoded
in Lean follows from its encoded hypotheses.  Source correspondence establishes
that the encoded theorem represents the mathematical claim made by Davis and
Kahan, including scalar field, dimension, operator domain, gap geometry, and
norm class.  We report these two checks separately.

For the sine-theta theorem, source correspondence requires the finite
interval/exterior case together with the ordered half-infinite separation cases
proved in the source.  For Section~4 the correspondence has a different form:
Proposition~4.4 cannot be represented as a theorem at its printed scope because
it is false.  A finite-dimensional counterexample isolates the invalid
inequality in the source proof.

The source-facing Davis--Kahan layer includes real and complex Hilbert spaces,
compact and noncompact operator settings as appropriate, and unbounded
self-adjoint operators represented by linear maps with explicit domains.  The corresponding
infrastructure is factored into a reusable library layer under the
\texttt{TauCeti} namespace, separately from the Davis--Kahan source statements.

We use the Yu--Wang--Samworth (YWS) perturbation bounds
\citep{yu2015davis} as a secondary application.  Their population-gap form is
particularly useful in statistics: a fixed separation in the population
spectrum can be combined directly with a stochastic bound on matrix-estimation
error.  This finite-dimensional development reuses portions of the same
spectral and subspace geometry and gives an independent source-correspondence
check in a modern statistical setting.

\paragraph{Contributions.}
\begin{enumerate}
  \item We give a Lean treatment of all \DKResultCount{} Davis--Kahan targets
        selected by a source-based counting rule.  \DKProvedCount{} are proved
        at source scope; Proposition~4.4 is formally refuted.  Every target is
        build-verified and has an accepted source-correspondence review.
  \item We formalize the source-facing sine-theta theorem with the unbounded
        operator, real/complex, where-defined unitarily invariant norm, and full
        finite-or-half-infinite gap scope.  We make explicit the mathematical
        correspondence between the source's positive $\sin\Theta_0$ operator
        and the rectangular operator used by the Lean theorem.
  \item We give a machine-checked $\mathbb{R}^4$ trace-norm counterexample to
        Davis--Kahan Proposition~4.4, identify the invalid summation step in
        equation~(4.3), retain the valid neighboring majorization results, and
        prove a corrected full-displacement theorem for $Q$-norms.
  \item We develop the operator-theoretic foundations needed to retain the
        source's scope, including principal-angle geometry, approximation
        numbers and operator ideals, Borel spectral calculus, unbounded
        self-adjoint operators, direct rotations, and Sylvester separation.
  \item As a secondary finite-dimensional application, we formalize the
        theorem-level claims of YWS Theorems~1--3 and formally check two printed
        defects: a missing square in equation~(4) and an incorrect rank-boundary
        convention in Theorem~3.
\end{enumerate}

\section{Davis--Kahan scope and coverage}
\label{sec:coverage}

Coverage is counted at the level of results stated by Davis and Kahan.  We
include the four unnumbered Section~2 headline
theorems and every named theorem, proposition, lemma, and corollary established
in the paper.  Definitions, examples, open questions, proof-local claims, and
results explicitly attributed to other sources are not counted as independent
targets.  This produces \DKResultCount{} result-level targets.

\begin{table}[H]
\centering
\small
\begin{tabular}{lrrrl}
\toprule
Source location & Targets & Proved & Refuted & Main mathematical content \\
\midrule
Section 2 & 4 & 4 & 0 & $\sin\Theta$, $\tan\Theta$, and double-angle bounds \\
Section 3 & 8 & 8 & 0 & direct rotations and two-subspace geometry \\
Section 4 & 5 & 4 & 1 & extremality and unitarily invariant norms \\
Section 5 & 3 & 3 & 0 & Sylvester equations and singular-value limits \\
Section 6 & 7 & 7 & 0 & generalized sine/tangent theorems and pinching \\
Section 8 & 2 & 2 & 0 & branch selection and spectral repulsion \\
\midrule
Total & 29 & 28 & 1 & \\
\bottomrule
\end{tabular}
\caption{Result-level Davis--Kahan coverage.  ``Proved'' means a checked Lean
statement accepted by source-correspondence review at the target scope.
Proposition~4.4 remains in the denominator with disposition ``refuted''; its
formal record contains a checked counterexample and repaired statements.}
\label{tab:dk-coverage}
\end{table}

The complete result table is given in Appendix~\ref{app:dk-inventory}.  The
acceptance criterion has two parts.  A declaration must resolve in the normal
build, and a separate review must establish how the Lean statement corresponds
to the source.  When the representation is nonliteral, the correspondence must
also be justified mathematically.  The sine-theta norm representation in the
next section is the main example.

\section{The source-facing sine-theta theorem}
\label{sec:sine-theta}

\subsection{Trial and exact subspaces}

Davis and Kahan distinguish a reference operator $A$ from a perturbed operator
$A+H$.  In this section we follow the current Lean interface and write $A$ for
the perturbed full-space operator; thus the Lean parameter \texttt{A}
corresponds to the source's $A+H$.

Let $E$ be the ambient Hilbert space.  Let $A_0$ be a self-adjoint operator on a
trial-coordinate space $F$, and let
\[
  E_0:F\to E
\]
be an isometric coordinate map for the trial subspace.  The residual compares
applying the ambient operator after embedding with applying the trial operator
before embedding:
\begin{equation}
  R = AE_0-E_0A_0.
  \label{eq:residual}
\end{equation}
When the operators are unbounded, \cref{eq:residual} is asserted on
$\operatorname{dom}(A_0)$ and $R$ is the bounded residual that agrees there.
The Lean predicate \texttt{IsTrialResidual} packages the isometry, the required
domain inclusion, and this residual identity.

Let $F_0:H\to E$ be isometric coordinates for the desired exact spectral
subspace and let $F_1:G\to E$ be isometric coordinates for its orthogonal
complement.  A self-adjoint operator $\Lambda_1$ on $G$ describes the
complementary spectral block through
\[
  AF_1 = F_1\Lambda_1
\]
on its domain.  The predicate \texttt{IsExactSpectralDecomposition} also records
orthogonality and completeness of the two coordinate ranges.

\subsection{The sine operator and its rectangular representative}

The source angle operator acts on the trial-coordinate space.  If
$P=F_0F_0^*$ is the orthogonal projection onto the desired exact subspace, then
\[
  \cos\Theta_0 = (E_0^*PE_0)^{1/2},
  \qquad
  \sin\Theta_0 = (I-E_0^*PE_0)^{1/2}.
\]
The current Lean theorem places the following rectangular map directly inside
the operator ideal:
\begin{equation}
  S=(I-P)E_0.
  \label{eq:rectangular-sine}
\end{equation}
The two operators have different codomains:
$\sin\Theta_0:F\to F$ and $S:F\to E$.  Their relationship is
\[
  S^*S
  = E_0^*(I-P)E_0
  = I-E_0^*PE_0,
\]
so
\begin{equation}
  |S|=(S^*S)^{1/2}=\sin\Theta_0.
  \label{eq:modulus-sine}
\end{equation}
A polar decomposition $S=U|S|$ and the ideal laws imply that $S$ and $|S|$
have the same singular-value data and the same value under every supported
unitarily invariant norm.  Consequently
\[
  N(S)=N(\sin\Theta_0)
\]
whenever the norm is defined.  This modulus identity is the semantic bridge
between the source operator and the Lean conclusion.

\subsection{Gap geometry and the inequality}

The source includes three forms of spectral separation.  The familiar case
places one diagonal block in a finite interval $[\beta,\alpha]$ and the other
outside its open $\delta$-neighborhood.  The Appendix to Section~6 also permits
ordered half-infinite separation: one block may be semibounded above while the
other is semibounded below at distance at least $\delta$, in either order.  The
Lean predicate \texttt{FormBoundedSylvesterGap} has constructors for all three
configurations.  The source spectral hypotheses imply the corresponding
form-bounded separation used by the Sylvester estimate.

The source also writes a common-domain equality for the unbounded trial
configuration.  That equality discharges the forward domain condition consumed
by \texttt{IsTrialResidual}.  Source correspondence is checked from the combined
hypotheses, with this equality supplying the required domain condition.

With these definitions, the current source-facing result states the
where-defined inequality
\begin{equation}
  \delta\,N(S)\leq N(R),
  \qquad S=(I-F_0F_0^*)E_0,
  \label{eq:dk-sine}
\end{equation}
for real or complex separable Hilbert spaces and possibly unbounded
self-adjoint $A$, $A_0$, and $\Lambda_1$.  ``Where-defined'' means that the
operator-ideal gauge is finite on both displayed operators.  This matches the
source convention and leaves membership transfer for arbitrary symmetric
operator ideals as a separate theorem.

Appendix~\ref{app:lean-guide} summarizes the general Lean notation used in the
displayed signatures.

\begin{lstlisting}[style=leanpaper,
  caption={Lean signature for the current source-facing theorem, abbreviated for
  readability.  Here \texttt{S} denotes the inlined rectangular map
  $(I-F_0F_0^*)E_0$; the exact Lean source uses the same theorem name and
  expands \texttt{S} in both membership and gauge expressions.},
  label={lst:dk-sintheta-headline}]
theorem sinTheta_unbounded_formGap_whereDefinedUIN_rclike
    [TopologicalSpace.SeparableSpace E]
    (N : NormalizedSymmetricOperatorIdealFamily K)
    (A : LinearPMap K E E)
    (A0 : LinearPMap K F F)
    (Lambda1 : LinearPMap K G G)
    (E0 : ContinuousLinearMap K F E)
    (F0 : ContinuousLinearMap K H E)
    (F1 : ContinuousLinearMap K G E)
    (R : ContinuousLinearMap K F E)
    (hA : IsSelfAdjoint A)
    (hA0 : IsSelfAdjoint A0)
    (hLambda1 : IsSelfAdjoint Lambda1)
    (htrial : IsTrialResidual A A0 E0 R)
    (hexact : IsExactSpectralDecomposition A Lambda1 F0 F1)
    (delta : Real) (hdelta : 0 < delta)
    (hgap : FormBoundedSylvesterGap A0 Lambda1 delta) :
    N.Mem S ->
    N.Mem R ->
      delta * N.gaugeReal S <= N.gaugeReal R := by
\end{lstlisting}

The source theorem contains all three separation configurations encoded by its
gap hypothesis.  The ordered half-infinite cases have distinct hypotheses from
the finite interval/exterior specialization, so source correspondence is checked
against the complete family of gap hypotheses.

\section{The remainder of the Davis--Kahan development}
\label{sec:dk-rest}

The sine theorem is only one of the four Section~2 inequalities and one of the
\DKResultCount{} result-level targets.  The rest of the development follows the
paper's mathematical progression.

\paragraph{Direct rotations and two-subspace geometry.}
Section~3 develops the canonical direct rotation between two subspaces,
including the acute and nonacute cases, square-root characterizations, a
classification theorem for pairs of subspaces, compact angle-eigenvalue
classification, commutation properties, and reversal symmetry.  The
formalization treats real and complex Hilbert spaces and separates reusable
two-projection and polar-decomposition facts from source-facing aliases.

\paragraph{Extremality and unitarily invariant norms.}
Section~4 compares the direct rotation with arbitrary unitaries that carry one
subspace onto another.  The valid results include pointwise singular-value
extremality for the restricted displacement, unitarily invariant norm
minimality of that restricted displacement, basis-angle square-sum
extremality, and majorization for the squared full displacement.  The source therefore has a correct majorization theorem for
$(I-W)^*(I-W)$.  Full-displacement minimality for $I-W$ requires a stronger
claim, which fails for the trace norm in \cref{sec:prop44}.

\paragraph{Sylvester theory.}
Section~5 requires an operator-level lower bound for the Sylvester map and a
semibounded self-adjoint form of the estimate.  The formalization also proves
the strong-cutoff convergence needed to pass singular-value information through
unbounded approximations.  This layer supplies the operator-level separation
mechanism used by the source-facing sine theorem, preserving the source spectral
gap hypotheses throughout the argument.

\paragraph{Generalized sine and tangent theorems.}
Section~6 weakens the isometric trial-map hypotheses and develops generalized
single-angle estimates, reflection pinching, pairwise-gap square-norm bounds,
tangent machinery, and a finite-rank near-maximizer leakage estimate.  The
formalization retains the source's distinction between the Section~2 isometric
trial coordinates and the lower-frame-bound hypotheses used by the generalized
results.

\paragraph{Branch selection.}
The two counted results in Section~8 formalize the spectral branch-selection
arguments that convert smallness and repulsion hypotheses into the acute branch
needed by the rotation theory.  These results use the unbounded spectral
calculus and cover cases beyond finite matrix decompositions.

This organization is reflected in the code: reusable analytic and geometric
facts are placed below the source-facing package, while
\texttt{DavisKahan/Sources/DavisKahan1970} records the correspondence to the
1970 paper.  Appendix~\ref{app:dk-inventory} lists all \DKResultCount{} targets.

\section{Proposition 4.4: refutation and repair}
\label{sec:prop44}

Davis--Kahan Proposition~4.4 asserts, in the real finite-dimensional case, that
when the largest principal angle is at most $\pi/3$, the direct rotation
minimizes every unitarily invariant norm of the full displacement $I-W$ among
unitaries carrying one subspace onto the other.  The statement is false.

\subsection{A four-dimensional counterexample}

Let $E=\mathbb{R}^4$ and
\[
  U=\operatorname{span}\{e_0,e_1\}.
\]
Define the orthogonal map $W=\tfrac12 H$ using
\[
H=
\begin{pmatrix}
 1&-1&-1&-1\\
 1& 1& 1&-1\\
-1&-1& 1&-1\\
 1&-1& 1& 1
\end{pmatrix},
\qquad
V=W(U).
\]
Both principal angles between $U$ and $V$ are $\pi/4$, so the pair satisfies
the printed $\pi/3$ hypothesis.  The singular values of $I-W$ are
$(\sqrt2,\sqrt2,0,0)$, hence
\[
  \|I-W\|_1=2\sqrt2.
\]
For the direct rotation $R$ from $U$ to $V$, all four singular values of $I-R$
are $\sqrt{2-\sqrt2}$, giving
\[
  \|I-R\|_1=4\sqrt{2-\sqrt2}>2\sqrt2.
\]
Thus an admissible competitor has smaller trace-norm displacement than the
direct rotation.  The construction, its acuteness and principal-angle bounds,
the singular-value calculations, and the strict inequality are all checked in
Lean.

\subsection{Where the printed proof fails}

The counterexample also pinpoints the invalid step in the proof.  Write
$K=I-W$, and let $\Omega_1$ and $\Omega_2$ be the two principal-plane
projections used by Davis and Kahan.  Their equation~(4.3) requires, on this
configuration,
\begin{equation}
  \|K\|_4
  \geq
  \|K\Omega_1\|_2+\|K\Omega_2\|_2,
  \label{eq:dk43-false}
\end{equation}
where $\|\cdot\|_j$ denotes the Ky Fan $j$-norm.  The witness gives
\[
  \|K\|_4=2\sqrt2,
  \qquad
  \|K\Omega_1\|_2=\|K\Omega_2\|_2=2,
\]
so \cref{eq:dk43-false} would require $2\sqrt2\ge4$.  Orthogonality of the two
domain planes does not force the images under $K$ to be orthogonal, and the
singular-value sums of the restrictions therefore need not add.

Proposition~4.3 remains valid.  It concerns the squared displacement
$(I-W)^*(I-W)$, for which the formalization proves the source majorization
statement and its unitarily invariant norm consequences.  The two neighboring
claims therefore have different dispositions: squared-displacement
majorization is valid, while arbitrary unitarily invariant norm minimality of
the full displacement fails.

\subsection{A Q-norm repair}

The formal development proves a natural correction.  Call a unitarily invariant
norm $N$ a $Q$-norm if there is a unitarily invariant norm $M$ such that
\begin{equation}
  N(X)^2=M(X^*X)
  \label{eq:qnorm-def}
\end{equation}
for every $X$.  Proposition~4.3 applied to $M$ immediately yields the valid
full-displacement estimate
\begin{equation}
  N(I-R)\leq N(I-W)
  \label{eq:qnorm-repair}
\end{equation}
for the direct rotation $R$.  The Lean theorem proves
\cref{eq:qnorm-repair} over every \texttt{RCLike} scalar field.  Its angle
assumption is the acuteness condition needed to construct the direct rotation;
the printed $\pi/3$ threshold is absent.  Conversely, the trace norm used above is formally shown not
to satisfy this $Q$-norm predicate on the witness space.

Earlier work on canonical unitary completion already records failures for
Schatten norms below the quadratic regime, including trace-norm phenomena
\citep{bolotnikov2001minimal}.  The present counterexample is inside the
Davis--Kahan short-angle hypothesis and directly refutes Proposition~4.4 as
printed.  Later literature has repeated the short-angle claim
\citep{qiu2006grassmann}; the formal result therefore serves as a correction to
a still-propagated statement.

\section{Yu--Wang--Samworth as a statistical specialization}
\label{sec:yws}

YWS gives a finite-dimensional statistical specialization of the same
subspace-perturbation geometry.  Its population-gap formulation connects a
fixed spectral separation in the population matrix directly to estimation
error, and its formalization provides a second source comparison with algebraic
and boundary-convention defects.

Let $\Sigma,\widehat\Sigma\in\mathbb{R}^{p\times p}$ be symmetric, with
population eigenvalues $\lambda_1\ge\cdots\ge\lambda_p$.  For
$1\le r\le s\le p$, let $d=s-r+1$, and let $V,\widehat V$ contain orthonormal
eigenvectors at indices $r,\ldots,s$.  With the population gap
\[
  \Delta=
  \min\{\lambda_{r-1}-\lambda_r,\lambda_s-\lambda_{s+1}\}>0,
\]
YWS Theorem~2 gives
\begin{equation}
  \|\sin\Theta(\widehat V,V)\|_{\mathrm F}
  \le
  \frac{2\min\{\sqrt d\,\|\widehat\Sigma-\Sigma\|_{\mathrm{op}},
                    \|\widehat\Sigma-\Sigma\|_{\mathrm F}\}}
       {\Delta},
  \label{eq:yws-sintheta}
\end{equation}
and an aligned-frame bound with an additional factor $\sqrt2$.  Theorem~3
applies the same principle to left and right singular subspaces through the Gram
operators.  The source audit also tracks the supporting equations and boundary
conventions needed to compare these theorem statements with the printed paper;
the detailed audit table is given in Appendix~\ref{app:yws-coverage}.

\subsection{Two printed defects}

The first defect is algebraic.  Equation~(4) in YWS omits a square.  For real
unit vectors $u,v$, the corrected identity is
\begin{equation}
(2\langle v,u\rangle)^2(1-\langle v,u\rangle^2)
=
\frac14\|v-u\|^2(2-\|v-u\|^2)^2(4-\|v-u\|^2).
\label{eq:yws-equation4-corrected}
\end{equation}
The formalization proves a numerical counterexample to the printed formula and
the corrected identity in general.

The second defect is a boundary convention in Theorem~3.  The printed
assignment
$\sigma^2_{\operatorname{rank}(A)+1}:=-\infty$ can make the denominator
infinite when the selected block reaches the rank boundary, forcing a zero
upper bound even when the compared singular subspaces differ.  The Gram-matrix
argument itself uses the ambient spectrum, where singular values beyond the
rank but within the ambient dimension are zero.  The corrected formal theorem
therefore extends the singular values by zero through the ambient dimension and
places the terminal convention after that dimension.  A $2\times2$ rank-one
example in the formalization witnesses the failure of the printed convention.

These are two source defects.  Each correction propagates to several formal
statements and supporting identities; the appendix records those affected
statements separately.

\section{Reusable formal foundations}
\label{sec:foundations}

Retaining the source scope required substantially more than a finite matrix
eigenvalue API.  The current repository separates reusable mathematics from the
paper-facing layer.  \texttt{ForTauCeti} contains declarations in their final
\texttt{TauCeti.*} namespaces and may depend only on Mathlib, the pinned Tau Ceti
library, and other \texttt{ForTauCeti} modules.  The Davis--Kahan layer builds
on this foundation for the source-specific statements, examples, and audit
checks, keeping those paper-specific objects out of the reusable operator-theory
API.

The main foundation groups used by the formalization are as follows.

\begin{table}[H]
\centering
\small
\begin{tabularx}{\linewidth}{>{\raggedright\arraybackslash}p{0.27\linewidth}X}
\toprule
Area & Role in the Davis--Kahan development \\
\midrule
Projection and principal-angle geometry & Orthogonal projections, directed sine maps, angle operators, principal-angle sequences, and two-subspace decompositions. \\
Approximation numbers and operator ideals & Infinite-dimensional singular-value surrogates, Ky Fan gauges, symmetric norming functions, Schatten-type families, and where-defined unitarily invariant norm comparisons. \\
Polar decomposition and direct rotations & Canonical partial isometries, acute/nonacute direct rotations, displacement formulas, and extremality statements. \\
Borel spectral calculus & Projection-valued measures, spectral projections, measurable functional calculus, and spectral restrictions used beyond the continuous functional calculus. \\
Unbounded self-adjoint operators & Domain-carrying linear partial maps, self-adjointness, semiboundedness, spectral restrictions, resolvents, and graph/domain arguments. \\
Sylvester equations & Finite, bounded, and unbounded separation estimates underlying sine and tangent perturbation bounds. \\
Real/complex transfer & Complexification and descent arguments needed to retain the source's real and complex statements without duplicating the entire analytic development. \\
Finite-dimensional spectral geometry & Matrix singular subspaces, Gram operators, aligned frames, and the explicit Proposition~4.4 counterexample. \\
\bottomrule
\end{tabularx}
\caption{Reusable mathematical layers required by the current formalization.}
\label{tab:foundations}
\end{table}

Mathlib supplies the general Lean foundation and continuous functional calculus
\citep{mathlib2020,dedecker2025cfc}.  Parts of the projection-valued-measure,
Borel-calculus, and self-adjoint-operator infrastructure were copied, adapted,
or generalized from the Spectra development \citep{bornemann2026spectra}.  The
current reusable layer is being prepared for upstream integration with Tau Ceti
\citep{tauceti2026}.  Appendix~\ref{app:source-review} describes how reuse and
source correspondence are recorded without encoding provenance or fidelity
claims into theorem names.

\section{Related work}
\label{sec:related}

Classical perturbation theory for invariant subspaces includes the
Davis--Kahan results, Wedin's singular-subspace bounds, and later matrix
perturbation treatments \citep{davis1970rotation,wedin1972perturbation,stewart1990matrix}.
For unbounded operators and operator angles, later Sylvester-equation approaches
provide closely related estimates and modern formulations
\citep{grubisic2005sylvester,grubisic2010tan}.  YWS reformulates the
finite-dimensional eigenvector problem around a population eigengap, making the
bound convenient for statistical analysis \citep{yu2015davis}; subsequent work
uses related spectral perturbation tools throughout high-dimensional statistics
\citep{cape2019two,chen2021spectral}.

The formalization builds on Lean~4 and Mathlib
\citep{demoura2021lean4,mathlib2020}.  Large source-driven formalizations such
as Flyspeck and the Liquid Tensor Experiment demonstrate that checking a
specific mathematical source can simultaneously require substantial reusable
library construction \citep{hales2017kepler,scholze2022liquid}.  Spectra is
particularly close technically because it develops operator and spectral
infrastructure in Lean that can be reused for parts of the present project
\citep{bornemann2026spectra}.

Recent work on autoformalization has made source correspondence a separate
research problem from proof search.  FormalAlign evaluates agreement between
informal and formal statements, ShadowBench probes semantic alignment through
auxiliary implications, and EconCSLib publishes source text together with Lean
statements and correspondence judgments
\citep{lu2025formalalign,han2026shadowbench,garg2026econcs}.  Our treatment of
Davis--Kahan uses the same conceptual separation in a manually reviewed
full-paper setting: successful compilation establishes the formal theorem, while
a separate source review establishes that the theorem represents the intended
claim.  The half-infinite sine-theta gap and the nonliteral
$S$--$\sin\Theta_0$ norm correspondence show why this second step is needed even
when the underlying proof is complete.

\section{Conclusion}
\label{sec:conclusion}

We give a Lean formal treatment of all \DKResultCount{} result-level targets in
Davis and Kahan's 1970 paper under an explicit source-based counting rule.
Twenty-eight targets are proved at the scope stated in the source.  The one
exception, Proposition~4.4, is false as printed; a checked counterexample,
proof diagnosis, and $Q$-norm repair replace an attempted proof of the false
claim.

The formalization retains the operator-theoretic scope of the paper.  The
current sine-theta endpoint handles
real/complex separable Hilbert spaces, unbounded self-adjoint operators,
finite and half-infinite gap configurations, and the source's where-defined
unitarily invariant norms.  Establishing the correspondence required an
explicit modulus argument because the Lean theorem and the source use operators
on different spaces to represent the same sine-angle norm.

The resulting principal-angle, operator-ideal, spectral-calculus, unbounded
operator, direct-rotation, and Sylvester theory is factored below the
paper-specific theorem layer.  The YWS formalization gives a secondary
finite-dimensional statistical application and independently demonstrates the
value of checking source identities and boundary conventions.  Across both
sources, kernel verification is paired with an explicit mathematical comparison
between the formal statement and the printed claim.

\section*{Acknowledgments}

This material is based upon work supported by the Defense Advanced Research
Project Agency (DARPA) under Contract No. HR001125CE017. Any opinions, findings
and conclusions or recommendations expressed in this material are those of the
author(s) and do not necessarily reflect the views of the Defense Advanced
Research Project Agency (DARPA).

\bibliographystyle{unsrtnat}
\bibliography{references}

\clearpage
\appendix

\section{Reading the Lean signatures}
\label{app:lean-guide}
\label{app:reading-lean}

A Lean theorem declaration has the form
\[
  \texttt{theorem name (parameters) (hypotheses) : conclusion := proof}.
\]
Parameters introduce the mathematical objects under discussion.  A hypothesis
such as \texttt{(hA : IsSelfAdjoint A)} supplies the corresponding fact;
bundled predicates such as \texttt{IsTrialResidual} and
\texttt{IsExactSpectralDecomposition} collect related assumptions.  Braces mark
arguments normally inferred by Lean, while square brackets request typeclass
instances supplying ambient algebraic, analytic, or topological structure.  The
proposition after the final colon is checked by Lean's kernel
\citep{demoura2021lean4}.

In \Cref{lst:dk-sintheta-headline}, \texttt{RCLike} is the Mathlib abstraction
covering $\mathbb{R}$ and $\mathbb{C}$.  A type
\texttt{ContinuousLinearMap K F E} is a continuous linear map.  A type
\texttt{LinearPMap K E E} is a linear partial map whose
domain is a submodule; this is the representation used for possibly unbounded
operators.  \texttt{N.Mem T} states that the ideal gauge associated with $N$ is
finite on $T$, and \texttt{N.gaugeReal T} is the corresponding real value.  The
suffix \texttt{Real} refers to the codomain of the gauge and not to the scalar
field.

\section{Davis--Kahan result inventory}
\label{app:dk-inventory}

\Cref{tab:dk-full-inventory} records the maintained result-level denominator.
Every row is currently verified in the ordinary Lean build and has accepted
semantic/source-correspondence review. The disposition column distinguishes a
proof of the source claim from the checked refutation of Proposition~4.4.

\begin{longtable}{>{\raggedright\arraybackslash}p{0.16\linewidth}>{\raggedright\arraybackslash}p{0.55\linewidth}>{\raggedright\arraybackslash}p{0.18\linewidth}}
\caption{The \DKResultCount{} Davis--Kahan result-level targets.}
\label{tab:dk-full-inventory}\\
\toprule
ID & Result & Disposition \\
\midrule
\endfirsthead
\toprule
ID & Result & Disposition \\
\midrule
\endhead
\\path{S2-sin-theta} & Single-angle sine theorem & proved at source scope \\
\\path{S2-tan-theta} & Single-angle tangent theorem & proved at source scope \\
\\path{S2-sin-two-theta} & Double-angle sine theorem & proved at source scope \\
\\path{S2-tan-two-theta} & Double-angle tangent theorem & proved at source scope \\
\\path{DK-3.1-prop} & Acute direct rotation existence and uniqueness & proved at source scope \\
\\path{DK-3.2-prop} & Nonacute existence criterion & proved at source scope \\
\\path{DK-3.3-prop} & Principal square-root characterization & proved at source scope \\
\\path{DK-3.4-prop} & Square as a direct rotation & proved at source scope \\
\\path{DK-3.1-thm} & Classification of pairs of subspaces & proved at source scope \\
\\path{DK-3.1-cor} & Compact classification by angle eigenvalues & proved at source scope \\
\\path{DK-3.5-prop} & Angle commutation and eigenspace geometry & proved at source scope \\
\\path{DK-3.2-cor} & Reversal symmetry & proved at source scope \\
\\path{DK-4.1-prop} & Pointwise and singular-value extremality of the direct rotation & proved at source scope \\
\\path{DK-4.1-cor} & Unitarily invariant norm minimality of restricted displacement & proved at source scope \\
\\path{DK-4.2-prop} & Basis-angle square-sum extremality & proved at source scope \\
\\path{DK-4.3-prop} & Squared-displacement unitarily invariant norm minimality & proved at source scope \\
\\path{DK-4.4-prop} & Full-displacement claim of Proposition~4.4 & refuted as printed; repaired separately \\
\\path{DK-5.1-thm} & Banach-space Sylvester lower bound & proved at source scope \\
\\path{DK-5.2-thm} & Semibounded self-adjoint Sylvester theorem & proved at source scope \\
\\path{DK-5.1-lem} & Strong-cutoff convergence of singular values & proved at source scope \\
\\path{DK-6.1-lem} & Direct-sum norm comparison and converse & proved at source scope \\
\\path{DK-6.2-lem} & Reflection-pinch contraction & proved at source scope \\
\\path{DK-6.1-prop} & Sine proof, ambient limitation, and symmetric sine theorem & proved at source scope \\
\\path{DK-6.1-thm} & Generalized sine theorem & proved at source scope \\
\\path{DK-6.2-thm} & Pairwise-gap square-norm sine theorem & proved at source scope \\
\\path{DK-6.3-thm} & Generalized tangent theorem and supporting machinery & proved at source scope \\
\\path{DK-6.3-lem} & Finite-rank near-maximizer leakage estimate & proved at source scope \\
\\path{DK-8.1-thm} & Branch selection and spectral repulsion & proved at source scope \\
\\path{DK-8.2-thm} & Smallness selects the acute branch & proved at source scope \\
\bottomrule
\end{longtable}

The result table is keyed to source results.  One source result can require many
Lean declarations, while a single reusable Lean theorem can support more than
one source-facing wrapper.  The denominator therefore remains stable under
refactoring of implementation modules.

\section{Source-correspondence review}
\label{app:source-review}

A result is accepted only after two distinct checks.

\paragraph{Mechanical verification.}
The declaration must resolve in the ordinary project build, and every theorem
on which it depends must be accepted by Lean's kernel.  This establishes the
formal proposition exactly as encoded.

\paragraph{Source comparison.}
The encoded proposition is compared with the cited source location.  The review
checks the mathematical objects, hypotheses, conclusion, scalar field,
dimension, norm class, boundary conventions, and operator-domain assumptions.
For a nonliteral representation, the review also records the theorem or
mathematical argument that justifies the correspondence.  Equation
\cref{eq:modulus-sine} is the main example: the rectangular map in the Lean
sine-theta theorem represents the source's positive sine operator only after a
modulus/polar-decomposition argument.

This separation prevents theorem names from serving as evidence of fidelity.
For example, the proved declaration
\texttt{sinTheta\_unbounded\_intervalExterior\_characterizedWitness\_rclike}
remains a useful finite-gap theorem, but it is not the selected source endpoint
because it does not include the half-infinite gap cases.  Conversely,
Proposition~4.4 is complete in the inventory even though its printed assertion
has no proof: the completed formal treatment consists of the transcribed claim
as a proposition, a checked counterexample/refutation, a diagnosis of the
source proof, and repaired theorems.

\section{YWS source-audit scope}
\label{app:yws-coverage}

The YWS audit uses a different denominator because the journal paper does not
claim a line-by-line formalization of every mathematical statement in that
article.  The maintained census has \YWSCensusEntryCount{} rows:
\YWSCensusPaperSourceCount{} rows associated with printed source material and
\YWSCensusExplicitAdditionCount{} explicitly marked additions.  All
\YWSCensusVerifiedCount{} rows currently resolve to declarations in the normal
build.

The source rows include the conclusions of Theorems~1--3 and the equations,
lemmas, examples, and boundary conventions needed to check those conclusions.
Several rows can share one correction.  In particular, the current census has
\YWSCensusCorrectedCount{} rows whose implementation is classified as corrected,
but the manuscript reports two printed YWS defects: the missing square in
equation~(4) and the rank-boundary convention in Theorem~3.

\section{Formalization procedure and AI assistance}
\label{app:formalization-process}

The development used a recurring five-stage workflow.
\begin{enumerate}
  \item \textbf{Gather context.}  Read the source result and surrounding proof,
        identify standing assumptions and notation, and inspect the relevant
        Lean interfaces.
  \item \textbf{Decompose.}  Express the target as mathematical obligations
        small enough to compare individually with existing library facts.
  \item \textbf{Find foundations.}  Search Mathlib, Tau Ceti, the local reusable
        layer, other Lean developments, and the mathematical literature before
        introducing new theory.
  \item \textbf{Formalize.}  Implement definitions and lemmas, compile them with
        Lean, and revise failures until the desired theorem is kernel-checked.
  \item \textbf{Review skeptically.}  Compare the resulting statement with the
        source independently of successful compilation.  Reopen the result if
        the review finds a scope, convention, representation, or hypothesis
        mismatch.
\end{enumerate}
The loop terminates for a source target only when both the formal proof and the
source-correspondence review are accepted, or when the source statement is
shown false and the formal record contains the corresponding refutation.

Large language models were used for proof planning, repository search, Lean
implementation, proof repair, and source comparison.  Human review determined
the intended paper scope, evaluated mathematical correspondence, and accepted
or rejected proposed repairs.  All formal propositions are certified by Lean's
kernel; model outputs supplied proof text and analysis that remained subject to
compilation and review.

\section{Artifact organization}
\label{app:artifact}

The paper-facing Davis--Kahan declarations are under
\path{DavisKahan/Sources/DavisKahan1970}.  Reusable mathematical infrastructure
is primarily under \path{ForTauCeti}; its declarations already use final
\texttt{TauCeti.*} namespaces.  The maintained Davis--Kahan result inventory is
\path{dev/davis-kahan-1970-formalization-result-inventory.json}.  The YWS source
census is maintained separately under the YWS package and is summarized for
this manuscript by the generated census snapshot.  This organization permits
source-audit records and paper-specific counterexamples to remain local to the
paper package while reusable analysis is prepared for upstream integration.


\end{document}
